\section{Preview and Metadata}
\label{sec:preview}

\subsection{The preview pane}

\begin{req}{PRV-1}
Selecting an entry shows a preview of it in a pane beside the list, together with
its details (\rid{META-1}) and extended attributes (\rid{META-2}).  The pane is shown
or hidden with the option \emph{Preview pane} or the Space key, and its width, between
240 and 900 pixels, is changed by dragging its left edge or from the keyboard.  It
starts open only on a window wide enough to hold it beside the list.
\end{req}

\begin{req}{PRV-2}
Which preview is shown is decided from the \emph{contents} of the file, never from its
name alone.  The name is only a hint about which check to try first: a file named
\code{.png} that is not an image is shown as ordinary text or as a binary file, and a
file with another name whose bytes are a PDF is a PDF.
\end{req}

\begin{req}{PRV-3}
Text is shown from the first 64~KB of the file and is always displayed as text, never
interpreted as markup.  When the file is longer, the pane says how much is shown.  A
file is text only if that part is valid UTF-8 without control characters; a binary
file is reported as such, with its size, and its bytes are not sent to the browser.  An
empty file is reported as empty.
\end{req}

\begin{req}{PRV-4}
Source code is shown with syntax highlighting, chosen from the name or extension, for
these languages: Bash, C, C++, CSS, diff, Dockerfile, Go, INI and TOML, Java,
JavaScript, JSON, Kotlin, Makefile, Markdown, PHP, Python, Ruby, Rust, SQL, Swift,
TypeScript, XML and HTML, and YAML.
\end{req}

\begin{req}{PRV-5}
A complete JSON file is shown pretty-printed; a JSON file that was cut off by
\rid{PRV-3} is shown as highlighted source.  A CSV or TSV file is shown as a table of
its first 200 rows, and the pane says when there are more.
\end{req}

\begin{req}{PRV-6}
A file whose leading bytes identify it as PNG, JPEG, GIF or WebP is shown as an
image.  Files of any other image type are not rendered: in particular SVG and HTML are
never rendered, because they can carry script.
\end{req}

\begin{req}{PRV-7}
A Markdown file (name ending \code{.md} or \code{.markdown}) is shown rendered by
default, with a \emph{Source} toggle to see the text.  Rendering supports headings,
emphasis, lists, task lists, tables, code, block quotes and rules.  Raw HTML in the
file is shown as text and never interpreted.  A link to another file opens that file
inside \texttt{fsb}, in its folder with the file selected; a link to a web address
opens in a new browser tab without sending a referrer; a link of any other kind
(\code{javascript:}, \code{data:}, \code{file:} and so on) is shown as plain text.  An
image that lies inside the roots is shown; a remote image is replaced by a note that
it was blocked, so previewing a file never contacts anything outside the machine.  At
most 20 images of at most 5~MB each are loaded for one file.
\end{req}

\begin{req}{PRV-8}
A file that begins with \code{\%PDF-} is shown in the browser's built-in PDF viewer.
A file named \code{.pdf} that does not begin so is shown as ordinary text or as a
binary file.
\end{req}

\begin{req}{PRV-9}
A file recognized by its bytes as a zip, tar or gzip-compressed tar archive (including
\code{.jar}) is shown as a table of its contents: each member's name and size, the
format, and the number of entries.  Nothing is extracted.  At most 1{,}000 entries are
shown; the listing reports when it stopped at a safety limit
(Section~\ref{sec:limits}).  Member names are displayed as text with control
characters replaced and are never used as paths.  A file that has an archive name but
is not one is shown as ordinary text or as a binary file.
\end{req}

\begin{req}{PRV-10}
For a file stored only in the cloud, every preview reports that it is not downloaded
and reads nothing (\rid{SEC-13}).
\end{req}

\begin{req}{PRV-11}
A symbolic link is previewed as what it points to, and the pane names the target.
\end{req}

\subsection{Hover peek}

\begin{req}{HOV-1}
Holding the pointer over a readable file for about 400~ms shows a small bubble with
its first 20 or so lines, as plain text.  Binary files, folders, empty files and
cloud-only files produce no bubble.
\end{req}

\begin{req}{HOV-2}
The bubble is dismissed by moving the pointer away, scrolling, or pressing a key.
Results are remembered so that hovering again does not re-read the file, and the
number of simultaneous requests is limited.  The option \emph{Hover previews} turns
the feature off; the choice is remembered (Section~\ref{sec:prefs}).
\end{req}

\begin{req}{HOV-3}
Hover peek passes through the same rules as every other access.  Because hovering
reveals contents more readily than choosing, files that commonly hold secrets are
denied by default (Section~\ref{sec:core}) and therefore never produce a bubble.
\end{req}

\subsection{Details and extended attributes}

\begin{req}{META-1}
The pane lists the entry's path, kind, size (in a readable unit and in bytes),
modification time and permissions, the target of a symbolic link, and whether the
file is stored only in the cloud.
\end{req}

\begin{req}{META-2}
Extended attributes are listed with their values.  A value is shown as text if it is
valid printable UTF-8 and as hexadecimal otherwise (labelled \emph{hex}, the label
before the digits).  At most 64 attributes are shown, and a value larger than
4~KB is shown by name and size only.  Attributes are available on macOS and Linux;
elsewhere none are shown.
\end{req}

\begin{req}{META-3}
The optional \emph{Attributes} column, hidden by default and placed last, shows
attribute \emph{names} only, for example \code{com.apple.quarantine}.  It is fetched
only for the rows on screen and remembered, and is never part of the listing itself,
so it cannot slow a large folder.
\end{req}
